\chapter{Fundamentação Teórica}
\label{chap_fundamentacao_teorica}

Autômatos celulares são sistemas dinâmicos discretos em espaço, tempo e estado. Em uma malha com \(L\) sítios, a configuração no instante \(t\) é
\[
\mathbf{x}^{(t)} = \big(x_0^{(t)},x_1^{(t)},\dots,x_{L-1}^{(t)}\big),
\qquad x_i^{(t)}\in\{0,1\}.
\]

Esse formalismo foi consolidado em estudos clássicos sobre auto-organização, modelagem discreta e computação em sistemas complexos \cite{vonneumann1966,toffoli1987,chopard1998,ilachinski2001}.

No caso dos autômatos celulares elementares (ECA), a regra local tem raio \(r=1\) e depende da vizinhança \((x_{i-1}^{(t)},x_i^{(t)},x_{i+1}^{(t)})\):
\[
x_i^{(t+1)} = f\big(x_{i-1}^{(t)},x_i^{(t)},x_{i+1}^{(t)}\big),
\]
com contorno periódico (índices módulo \(L\)). Como
\[
f:\{0,1\}^3\to\{0,1\},
\]
existem \(2^8=256\) regras possíveis, numeradas de 0 a 255 \cite{wolfram1983,wolfram2002,mathworldECA}. A dinâmica global pode ser escrita como
\[
\mathbf{x}^{(t+1)} = F\big(\mathbf{x}^{(t)}\big),
\qquad F:\{0,1\}^L\to\{0,1\}^L.
\]

Na visualização espaço-temporal, adota-se a convenção usual de ECA: \(t=0\) na linha superior e tempo crescente para baixo \cite{wikipediaECA,mathworldECA}.

Na codificação de Wolfram, para \(a,b,c\in\{0,1\}\), define-se
\[
p = 4a + 2b + c \in \{0,\dots,7\}.
\]
Se
\[
R = \sum_{k=0}^{7} b_k 2^k,
\qquad b_k\in\{0,1\},
\]
então, na ordem \(111,110,\dots,000\), a atualização local é
\[
x_i^{(t+1)} = b_{7-p}.
\]

\section{Métricas de análise}

As métricas usadas para interpretar os padrões são consistentes com a análise de transições de fase e complexidade em autômatos celulares \cite{langton1990,li1990,wolfram2002}:

\textbf{Densidade de sítios ativos}
\[
\rho(t)=\frac{1}{L}\sum_{i=0}^{L-1} x_i^{(t)}.
\]

\textbf{Atividade temporal}
\[
A(t)=\frac{1}{L}\sum_{i=0}^{L-1}\left|x_i^{(t+1)}-x_i^{(t)}\right|.
\]

\textbf{Entropia binária instantânea}
\[
H(t)=-\rho(t)\log_2\rho(t)-\big(1-\rho(t)\big)\log_2\big(1-\rho(t)\big),
\]
adotando \(0\log 0=0\).

\textbf{Correlação espacial de alcance \(r\)}
\[
C_t(r)=\frac{1}{L}\sum_{i=0}^{L-1}x_i^{(t)}x_{i+r}^{(t)}-\rho(t)^2.
\]

\textbf{Espalhamento de dano}
\[
d_H(t)=\frac{1}{L}\sum_{i=0}^{L-1}\left|x_i^{(t)}-y_i^{(t)}\right|,
\]
com \(\mathbf{x}^{(t)}\) e \(\mathbf{y}^{(t)}\) evoluindo pela mesma regra a partir de condições iniciais próximas.

\section{Classes de Wolfram}

A classificação clássica de Wolfram separa os ECA em quatro regimes \cite{wolfram1983,wolfram2002}:

\textbf{Classe 1 (Homogênea)}: convergência para ponto fixo,
\[
\exists T:\forall t\ge T,\; x^{(t)}=x^*, \qquad F(x^*)=x^*.
\]

\textbf{Classe 2 (Periódica)}: convergência para ciclo limite,
\[
\exists T,\,p\ge 1:\forall t\ge T,\; x^{(t+p)}=x^{(t)}.
\]

\textbf{Classe 3 (Caótica)}: alta sensibilidade a perturbações,
\[
\|F^t(x)-F^t(x+\delta)\|\not\to 0.
\]

\textbf{Classe 4 (Complexa)}: coexistência de ordem local e estruturas propagantes (gliders), com interações não triviais.

No contexto de computação, a Regra 110 (Classe 4) é universal no sentido de Turing \cite{cook2004}. Assim, regras locais simples podem produzir dinâmicas de alta complexidade.

As definições acima são consistentes com a teoria clássica e fundamentam a análise qualitativa apresentada neste trabalho \cite{wolfram1983,wolfram2002,langton1990,li1990}.